\section{Compresión}
\label{sec:compresion}

La compresión es el proceso de reducir el tamaño de los datos para disminuir el espacio de almacenamiento requerido y el tiempo de transmisión. Se trata de una técnica fundamental en el procesamiento digital, con aplicaciones como la compresión de imágenes (por ejemplo, \texttt{JPEG}, \texttt{PNG}) y de audio (por ejemplo, \texttt{MP3}, \texttt{FLAC}).

En función del tratamiento de la información, existen dos tipos principales de compresión \cite{Kavitha2016A}:

\begin{itemize}
    \item Compresión sin pérdida (\textit{lossless}): permite reconstruir los datos originales de forma exacta a partir de los datos comprimidos, sin pérdida de información.
    \item Compresión con pérdida (\textit{lossy}): sacrifica parte de la información considerada menos perceptible o relevante para lograr una tasa de compresión mayor. En este caso, se busca que la diferencia entre el dato original y el reconstruido resulte imperceptible desde un punto de vista perceptual.
\end{itemize}

Este último enfoque resulta especialmente relevante en la compresión de redes neuronales, donde se sacrifica precisión numérica en favor de la eficiencia computacional.
